\section{Related literature}
\label{sec:literature}

The model sits in three literatures: the tax--benefit microsimulation tradition that begins with \citet{orcutt1957}, the survey-data literature on under-reporting and calibration that determines how far a population estimate can be trusted, and the open-source and reproducibility movement that determines whether anyone outside the producing institution can check it.

\subsection{Tax--benefit microsimulation}

\citet{orcutt1957} proposed simulating a population as a collection of individual units whose behaviour and circumstances are represented explicitly, rather than as aggregate relationships between averages; the tax--benefit branch of that programme applies statutory rules to a representative sample of households and aggregates the results. The survey and handbook literature---\citet{mitton2000}, \citet{bourguignon2006}, \citet{figari2015}---sets out the standard architecture: a microdata file of households, a rule engine encoding the tax and transfer system for a given policy year, and a differencing procedure that runs the same data under baseline and reform parameters and reports the change in revenue, in net incomes, and in their distribution.

Every major fiscal institution operates a model of this class. In the United Kingdom, the Institute for Fiscal Studies' TAXBEN has been the standard academic instrument for four decades, and EUROMOD extended the design to a harmonised multi-country platform whose openness of documentation---if not of data---set the transparency benchmark for the field \citep{sutherland2013}. In the United States, the NBER's TAXSIM has provided federal and state tax calculations to researchers since the 1970s \citep{feenberg1993}, and the Policy Simulation Library's Tax-Calculator provides an openly developed federal model \citep{taxcalculator}. The revenue departments and fiscal watchdogs run their own: HMRC's costings, from which the ready reckoner used in Section~\ref{sec:validation} is produced, rest on administrative Survey of Personal Incomes data rather than a household survey \citep{hmrc2025reckoner}.

PolicyEngine's distinguishing choices within this tradition are the ones that matter for this paper. The rule engine is open source and so is the model of statute: \texttt{policyengine-core} is a fork of the OpenFisca framework \citep{openfisca}, in which every quantity---an allowance, a taper, a means test---is a \emph{variable} with a formula and a set of dated \emph{parameters}, and a reform is a re-dating of a parameter rather than an edit to code. And the same object serves both levels: the identical reform dictionary that scores one typed household scores the whole survey, which is why the exactness question of Section~\ref{sec:uncertainty} arises so sharply---the arithmetic is literally the same code in both cases, and only the inputs differ in status.

\subsection{Survey microdata, under-reporting, and calibration}

The credibility of a population estimate rests on the microdata, and the microdata literature is not reassuring about using a household survey raw. \citet{meyer2015} document a broad and worsening under-reporting of transfer receipt in US household surveys, with major programmes capturing well under the administratively recorded totals; \citet{brewer2017} show for the UK Family Resources Survey that the households reporting the lowest incomes are conspicuously better off on expenditure measures than their reported incomes imply, so the bottom of the reported income distribution is substantially a measurement artefact. Both findings bear directly on a distributional costing: the decile a household is assigned to, and the benefit entitlement computed for it, depend on income components the survey measures with error.

The standard responses are imputation and calibration. Reweighting a survey so that its weighted totals reproduce known administrative or national-accounts aggregates is the calibration-estimator problem of \citet{deville1992}, and static ageing---uprating incomes and re-weighting a survey of one year to represent a later policy year---is routine in the microsimulation literature \citep{figari2015}. PolicyEngine's UK data pipeline does both: the \emph{enhanced} Family Resources Survey \citep{dwpfrs} imputes under-reported incomes and reweights the survey so aggregates better match administrative totals, and the per-year datasets used for population runs are built by uprating to the policy year. This is what makes population estimates usable; it is also, precisely, what makes them estimates. A reweighted, imputed, aged survey is a model of the population, and its errors are not the sampling error of the original design alone.

\subsection{Open models and reproducibility}

The third literature is about who can check the answer. Systematic replication attempts find that published economic results frequently cannot be reproduced from the authors' own materials \citep{chang2015}, and the best-known failure in the fiscal domain---the spreadsheet error in Reinhart and Rogoff's debt-threshold result, found only when the original file was obtained and re-run \citep{herndon2014}---is a standing argument that consequential numbers should be recomputable by outsiders. Official costings are the strongest case: they move budgets, and the models producing them are typically not runnable outside the producing department.

An open microsimulation answers part of that objection and, honestly, only part. The rules are public code and the household calculation is reproducible by anyone with a laptop and no credentials at all---which is a genuinely unusual property in fiscal analysis, and the reason Section~\ref{sec:validation}'s worked example can be checked by hand. The population side is not fully open in the same way: the UK microdata is distributed through a gated repository requiring a token, so a reader can reproduce this paper's household arithmetic exactly and cannot reproduce its \pounds 6.46bn costing without being granted access. Section~\ref{sec:repro} states that asymmetry rather than eliding it. The same tension---rules open, microdata restricted---runs through the whole field, since the underlying surveys are released under licence by national statistical agencies.

Finally, the multi-model design this paper's model leads is itself a position in the modelling literature: rather than extend one model to answer every question, the suite scores the same reform through models of different classes and publishes a \emph{verification gradient} recording how far each one's outputs can be audited against ground truth, in the tradition of institutional model suites \citep{burgess2013}. The microsimulation occupies the top of that gradient---``checked against implemented legislation''---and the sister papers \citep{ahmadi2026obr, ahmadi2026og} document the members below it.
